\documentclass[11pt]{article}

% Text
\usepackage[ngerman]{babel}
\usepackage[utf8]{inputenc}
\usepackage[official]{eurosym}
\usepackage{enumitem}
\usepackage{caption}
    \captionsetup{format=hang}
%\usepackage[center]{titlesec}
%\setcounter{secnumdepth}{0}
\usepackage[dvipsnames]{xcolor}
    \colorlet{darkgreen}{green!40!black}
    \colorlet{darkyellow}{yellow!87!black}
    \colorlet{DarkBrown}{Brown!60!white}
    \colorlet{YellowBrown}{Brown!40!Yellow}
    \colorlet{darkblue}{blue!60!black}
    \colorlet{whiteyellow}{darkyellow!60!white}
    
    \colorlet{mylightblue}{blue!20}
    \colorlet{myblue}{blue!50!black}
    \colorlet{mydarkblue}{blue!30!black}
    \colorlet{mylightred}{red!10}
    \colorlet{myred}{red!50!black}
    \colorlet{mydarkred}{red!60!black}
    \colorlet{mydarkgreen}{green!30!black}
\usepackage{contour}
\usepackage{hyperref}
\hypersetup{colorlinks=true,
            linkcolor=black,
            urlcolor=black}
\parindent0cm
%\pagestyle{empty}

% Maths
\usepackage{amssymb, amsmath, mathtools, siunitx, bm}
\usepackage{graphicx}
\usepackage[a4paper, left=2cm, right=2cm, top=2cm]{geometry}
\usepackage{makecell}
\usepackage{vcell}
\usepackage{booktabs}
\usepackage{float}
%    \restylefloat{table}
\usepackage{pgfplots, pgfplotstable, filecontents}
\pgfplotsset{compat=1.18, width=10cm}
    \newcommand{\linregAddplot}[2][]{%
      % Will plot a linear regression from a datafile.
      %
      % Usage: \linregAddplot[<options>]{<file>}
      % where <options> is the same options as for \addplot+ and <file> is a simple
      % datafile gnuplot can read.
      %
      % The macros \linregSlope, \linregIntercept and \linregRsq are available after
      % \linregAddplot has been called. They can be used in
      % \addlegendentry{}.
    
        \addplot+[raw gnuplot, no markers, #1] gnuplot {
        stats '#2' nooutput; % Calculate a bunch of statistical values
        plot [x=STATS_min_x:STATS_max_x] STATS_slope*x + STATS_intercept;
        set print "#2.linreg";
        print STATS_slope; % Write the parameters to file
        print STATS_intercept;
        print STATS_correlation;
      };
      \IfFileExists{./#2.linreg}{
        \expandafter\newread\csname read#2\endcsname % Open the file again
        \expandafter\openin\csname read#2\endcsname=#2.linreg
    
        % Read the values from the file
        \expandafter\read\csname read#2\endcsname to \mytmp
        % and pprint them to tmp macros
        \pgfmathprintnumberto{\mytmp}{\linregSlope}
    
        \expandafter\read\csname read#2\endcsname to \mytmp
        \pgfmathprintnumberto{\mytmp}{\linregIntercept}
    
        \expandafter\read\csname read#2\endcsname to \mytmp
        % Square the r value before pprinting
        \pgfmathpow{\mytmp}{2}
        \pgfmathprintnumberto[precision=5]{\pgfmathresult}{\linregRsq}
          }{%
        \def\linregSlope{NaN}
        \def\linregIntercept{NaN}
        \def\linregRsq{NaN}
        \message{You need to run LaTeX again to get you linear regression labels right.}
      }
    }

% Chemistry
\usepackage{modiagram}
\usepackage{chemfig}
\usepackage[version=4]{mhchem}
\usepackage{chemformula}
\usepackage{elements}
\usepackage{isotope}

% Grafics
\usepackage{tikz}
\usetikzlibrary{shapes.geometric,
                arrows,
                positioning,
                hobby,
                arrows.meta,
                calc,
                intersections,
                decorations.markings}
\usepackage{ifthen}

% Born-Haber Kreisprozess
\tikzset{box/.style={rectangle, minimum size=5mm, draw=black, fill=black!10}}
\tikzset{arrow/.style={->, very thick, >=stealth}}

% Phasendiagramm
\tikzset{>={Latex[length=4,width=4]}}
\tikzset{midarr/.style={decoration={markings,
                                    mark=at position #1 with {\arrow{stealth}}},
                                    postaction={decorate}},
         midarr/.default=0.5}
\def\tick#1#2{\draw[thick] (#1) ++ (#2:0.03*\ymax) --++ (#2-180:0.06*\ymax)}
\newcommand\namebond[5]
                    [-1pt]
                    {\chemmove{\path(#2)--(#3)node[midway,
                                                   #4,
                                                   yshift=#1,
                                                   black!60]
                                                  {#5};}}
\newcommand\arcbetweennodes[3]
                           {\pgfmathanglebetweenpoints{\pgfpointanchor{#1}{center}}
                                                      {\pgfpointanchor{#2}{center}}
                            \let#3\pgfmathresult}
\newcommand\arclabel[6]
                    [red,
                     -stealth,
                     shorten <=1pt,
                     shorten >=1pt]
                    {\chemmove{\arcbetweennodes{#4}{#3}\anglestart
                               \arcbetweennodes{#4}{#5}\angleend
                               \ifdim\anglestart pt>\angleend pt \pgfmathsetmacro\anglestart{\anglestart-360}\fi
                               \draw[#1]
                                    ([shift=(\anglestart:#2)]#4)arc
                                    [start angle=\anglestart,
                                     end angle=\angleend,
                                     radius=#2];
                               \pgfmathsetmacro\anglestart{(\anglestart+\angleend)/2}
                               \node[shift=(\anglestart:#2+1pt)#4,
                                     anchor=\anglestart+180,
                                     inner sep=0pt,
                                     outer sep=0pt]at
                                    (#4)
                                    {#6};}}
                                    
% Orbitale
\pgfdeclarelayer{background}
\pgfdeclarelayer{foreground}

\pgfsetlayers{background,main,foreground}

\colorlet{electron}{black!100}                                         % Elektronenfarbe

\tikzset{px/.style={thick,fill opacity=1}}
\tikzset{px 0/.style={px,fill=white}}                                  % px-Orbitalfarbe (0 Elektronen)
\tikzset{px 1/.style={px,fill=white}}                                  % px-Orbitalfarbe (1 Elektron)
\tikzset{px 2/.style={px,fill=white}}                                  % px-Orbitalfarbe (2 Elektronen)

\tikzset{py/.style={thick,fill opacity=.5}}
\tikzset{py 0/.style={py}}                                             % py-Orbitalfarbe (0 Elektronen)
\tikzset{py 1/.style={py}}                                             % py-Orbitalfarbe (1 Elektron)
\tikzset{py 2/.style={py}}                                             % py-Orbitalfarbe (2 Elektron)

\tikzset{pz/.style={thick,fill opacity=1}}
\tikzset{pz 0/.style={pz,fill=white}}                                  % pz-Orbitalfarbe (0 Elektronen)
\tikzset{pz 1/.style={pz,fill=white}}                                  % pz-Orbitalfarbe (1 Elektron)
\tikzset{pz 2/.style={pz,fill=white}}                                  % pz-Orbitalfarbe (2 Elektron)

\tikzset{spx/.style={thick,fill opacity=1}}
\tikzset{spx 0/.style={spx,fill=white}}                                 % sp-Orbitalfarbe (0 Elektronen)
\tikzset{spx 1/.style={spx,fill=white}}                                 % sp-Orbitalfarbe (1 Elektron)
\tikzset{spx 2/.style={spx,fill=white}}                                 % sp-Orbitalfarbe (2 Elektron)

\tikzset{atomcore/.style={shape=circle,
                          thick,
                          draw=black,                                  % Atomkernfarbe
                          minimum size=7mm,                            % Atomkerngröße
                          font=\large,                                 % Atomkernschriftgröße
                          inner sep=0pt}}

\def\pxheight{1.75}                                                    % px-Orbitalhöhe
\def\pxwidth{.3}                                                       % px-Orbitalweite

\def\pyheight{1.625}                                                   % py-Orbitalhöhe
\def\pywidth{.4375}                                                    % py-Orbitalweite

\def\pzheight{1.625}                                                   % pz-Orbitalhöhe
\def\pzwidth{.4375}                                                    % pz-Orbitalweite

\def\spxheight{1.75}                                                   % psx-Orbitalhöhe
\def\spxwidth{.3}                                                      % psx-Orbitalweite

\newcommand{\px}
           [3]
           {\begin{scope}[rotate=#1,shift=(#2)]
                \coordinate (c1)   at (-\pxwidth    ,.6*\pxheight);
                \coordinate (c2)   at (-\pxwidth    ,\pxheight);
                \coordinate (c3)   at (\pxwidth     ,\pxheight);
                \coordinate (c4)   at (\pxwidth     ,.6*\pxheight);
                \coordinate (top)  at (0            ,\pxheight);
        
                \coordinate (e1)   at (.42*\pxwidth ,.52*\pxheight);
                \coordinate (e2)   at (-.5*\pxwidth ,.68*\pxheight);
                \coordinate (e3)   at (-.36*\pxwidth,.7*\pxheight);
                \coordinate (e4)   at (.56*\pxwidth ,.54*\pxheight);
            \end{scope}
            
            \begin{pgfonlayer}{background}
                \draw[px #3] (#2) .. controls (c1) and (c2) .. (top) ..
                                     controls (c3) and (c4) .. (#2);
            \end{pgfonlayer}
            
            \ifnum#3=1
                \draw[arrows = {-Stealth[left]}] (e1) -- (e2);
            \else
            \ifnum#3=2
                \draw[arrows = {-Stealth[left]}] (e1) -- (e2);
                \draw[arrows = {-Stealth[left]}] (e3) -- (e4);
            \fi
            \fi}
\newcommand{\py}
           [3]
           {\begin{scope}[rotate=#1,shift=(#2)]
                \coordinate (c1)  at (-\pywidth      ,.6*\pyheight);
                \coordinate (c2)  at (-\pywidth      ,\pyheight);
                \coordinate (c3)  at (\pywidth       ,\pyheight);
                \coordinate (c4)  at (\pywidth       ,.6*\pyheight);
                \coordinate (top) at (0              ,\pyheight);
        
                \coordinate (e1)   at (.525*\pywidth ,.6*\pyheight);
                \coordinate (e2)   at (-.475*\pywidth,.6*\pyheight);
                \coordinate (e3)   at (-.475*\pywidth,.63*\pyheight);
                \coordinate (e4)   at (.525*\pywidth ,.63*\pyheight);
            \end{scope}
            
            \begin{pgfonlayer}{background}
                \draw[py #3] (#2) .. controls (c1) and (c2) .. (top) ..
                                     controls (c3) and (c4) .. (#2);
            \end{pgfonlayer}
            
            \ifnum#3=1
                \draw[arrows = {-Stealth[left]}] (e1) -- (e2);
            \else
            \ifnum#3=2
                \draw[arrows = {-Stealth[left]}] (e1) -- (e2);
                \draw[arrows = {-Stealth[left]}] (e3) -- (e4);
            \fi
            \fi}
\newcommand{\pz}
           [3]
           {\begin{scope}[rotate=#1,shift=(#2)]
                \coordinate (c1)  at (-\pzwidth     ,.6*\pzheight);
                \coordinate (c2)  at (-\pzwidth     ,\pzheight);
                \coordinate (c3)  at (\pzwidth      ,\pzheight);
                \coordinate (c4)  at (\pzwidth      ,.6*\pzheight);
                \coordinate (top) at (0             ,\pzheight);
        
                \coordinate (e1)  at (-.075*\pzwidth,.45*\pzheight);
                \coordinate (e2)  at (-.075*\pzwidth,.7*\pzheight);
                \coordinate (e3)  at (.075*\pzwidth ,.7*\pzheight);
                \coordinate (e4)  at (.075*\pzwidth ,.45*\pzheight);
            \end{scope}
            
            \begin{pgfonlayer}{background}
                \draw[pz #3] (#2) .. controls (c1) and (c2) .. (top) ..
                                     controls (c3) and (c4) .. (#2);
            \end{pgfonlayer}
            
            \ifnum#3=1
                \draw[arrows = {-Stealth[left]}] (e1) -- (e2);
            \else
            \ifnum#3=2
                \draw[arrows = {-Stealth[left]}] (e1) -- (e2);
                \draw[arrows = {-Stealth[left]}] (e3) -- (e4);
            \fi
            \fi}
\newcommand{\spx}
           [3]
           {\begin{scope}[rotate=#1,shift=(#2)]
                \coordinate (c1)   at (-\spxwidth    ,.6*\spxheight);
                \coordinate (c2)   at (-\spxwidth    ,\spxheight);
                \coordinate (c3)   at (\spxwidth     ,\spxheight);
                \coordinate (c4)   at (\spxwidth     ,.6*\spxheight);
                \coordinate (top)  at (0            ,\spxheight);
        
                \coordinate (e1)   at (-.5*\spxwidth ,.6*\spxheight);
                \coordinate (e2)   at (.35*\spxwidth ,.75*\spxheight);
                \coordinate (e3)   at (.5*\spxwidth ,.72*\spxheight);
                \coordinate (e4)   at (-.4*\spxwidth ,.57*\spxheight);
            \end{scope}
            
            \begin{pgfonlayer}{background}
                \draw[spx #3] (#2) .. controls (c1) and (c2) .. (top) ..
                                     controls (c3) and (c4) .. (#2);
            \end{pgfonlayer}
            
            \ifnum#3=1
                \draw[arrows = {-Stealth[left]}] (e1) -- (e2);
            \else
            \ifnum#3=2
                \draw[arrows = {-Stealth[left]}] (e1) -- (e2);
                \draw[arrows = {-Stealth[left]}] (e3) -- (e4);
            \fi
            \fi}
\newcommand{\Atom}
           [3]
           [AtomNode]
           {\node[atomcore] (#1) {\ce{#2}}; \foreach \ang/\anchor/\n in {#3} {\px{\ang}{#1.\anchor}{\n}}}
\newcommand{\drawpy}
           [3]
           [AtomNode]
           {                                \foreach \ang/\anchor/\n in {#3} {\py{\ang}{#1.\anchor}{\n}}}
\newcommand{\drawpz}
           [3]
           [AtomNode]
           {                                \foreach \ang/\anchor/\n in {#3} {\pz{\ang}{#1.\anchor}{\n}}}
\newcommand{\drawspx}
           [3]
           [AtomNode]
           {                                \foreach \ang/\anchor/\n in {#3} {\spx{\ang}{#1.\anchor}{\n}}}
           
% Erste Seite
\title{Mathe}
\author{Tobias Demmel}
\date{\today}

\begin{document}

\maketitle

\tableofcontents

\newpage

\section{Plots}

\subsection{2D-Plot einer Funktion}

    \begin{figure}[!ht]
        \centering
        \scalebox{0.75}{
        \begin{tikzpicture}
            \begin{axis}[xmin=-2.5, xmax=2.5,
                         ymin=-5, ymax=5,
                         axis lines=middle,
                         xlabel=$x$, ylabel=$y$,
                         title={$\boldsymbol{x^2}$ und $\boldsymbol{1-x^2}$}]
                \addplot[color=red, dashed, mark=*, samples=50, domain=-1.75:1]{x^2};
                \addplot[color=blue, samples=50]{1-x^2};
            \end{axis}
        \end{tikzpicture}
        }
        \caption{2D-Plot}
    \end{figure}
    
\subsection{3D-Plot einer Funktion}

    \begin{figure}[!ht]
        \centering
        \scalebox{0.75}{
        \begin{tikzpicture}
            \begin{axis}[view={40}{30},
                         zlabel=$z$,
                         xlabel=$x$,
                         ylabel=$y$,]
                \addplot3[color=red,
                          mesh,
                          samples=25]
                         {x^2 - y^2};
            \end{axis}
        \end{tikzpicture}
        }
        \caption{3D-Plot}
    \end{figure}
    
    \begin{figure}[!ht]
        \centering
        \scalebox{.75}{
        \begin{tikzpicture}
            \begin{axis}[xmin=-3, xmax=3,
                         ymin=-3, ymax=3,
                         zmin=0, zmax=2,
                         view={320}{20},
                         zlabel=$z$, xlabel=$x$, ylabel=$y$,
                         unit vector ratio = 1.5 1.5 1]
                \addplot3[mesh,
                          surf,
                          shader=faceted,
                          samples=40,
                          domain=-3:3]
                         {(x^2+2*y^2)*exp(-x^2-y^2+1)};
            \end{axis}
        \end{tikzpicture}
        }
        \caption{Stauchung (z-Achse)}
    \end{figure}
    
\subsection{Vektorfeld}

    \begin{figure}[H]
        \centering
        \scalebox{.625}{
        \begin{tikzpicture}
            \begin{axis} [view = {0}{90},
                          xmin=-1, xmax=1,
                          ymin=-1, ymax=1,
                          unit vector ratio = 2 1]
                \addplot3[quiver = {u = {x},
                                    v = {y},
                                    scale arrows = 0.25},
                          -stealth,
                          domain = -3:3]
                         {0};
            \end{axis}
        \end{tikzpicture}
        }
        \caption{Vektorfeld}
    \end{figure}
    
\subsection{2D-Plot einer Funktion}

    \begin{filecontents}{data.dat}
        x    y    ex   ey   uy   ly
        2.92 4.69 0.2  0    4.69 4.69
        3.00 4.33 0.15 0.05 4.38 4.28
        3.10 3.91 0.1  0.1  4.01 3.81
        3.19 3.47 0.05 0.15 3.62 3.32
        3.30 2.94 0    0.2  3.14 2.74
    \end{filecontents}
    \begin{figure}[!ht]
        \centering
        \scalebox{.75}{
        \begin{tikzpicture}
            \begin{axis}[xmin=2.895, xmax=3.325,
                         ymin=2, ymax=5.25,
                         axis lines=middle,
                         x post scale = 1.4,
                         xlabel=$\frac{1}{T}\cdot10^{-3}/\frac{1}{K}$, ylabel=$\ln{(P/mbar)}$,
                         x label style={at={(axis description cs:.5,-.1)}, anchor=north},
                         y label style={at={(axis description cs:-.1,.5)}, rotate=90, anchor=south},
                         xmajorgrids=false,
                         ymajorgrids=false,
                         grid style=dashed,
                         legend pos=outer north east]
                %Messwerte der Abkühlkurve
                \addplot[only marks,
                         mark = +] table
                        {data.dat};
                %Eichgerade (blau)
                \addplot[densely dotted,
                         thick,
                         blue] table
                        [y={create col/linear regression={y=y}}]
                        {data.dat} node
                        [right,
                         pos=0.25]
                        {\scriptsize\qquad$\pgfmathprintnumber{\pgfplotstableregressiona} \cdot x \pgfmathprintnumber[print sign]{\pgfplotstableregressionb}$};
                %Fehlerbalken
                \addplot[error bars/.cd,
                         x dir=both,
                         x fixed =0.0125,
                         y dir=both,
                         y explicit] table
                        [y error = ey,
                         only marks]
                        {data.dat};
                %Obere Grenzgerade (rot)
                \addplot[red,
                         dotted] table
                        [y={create col/linear regression={y=uy}}]
                        {data.dat} node
                        [right,
                         pos=0.6]
                        {\scriptsize\quad$\pgfmathprintnumber{\pgfplotstableregressiona} \cdot x \pgfmathprintnumber[print sign]{\pgfplotstableregressionb}$};
                %Untere Grenzgerade (rot)
                \addplot[red,
                         dotted] table
                        [y={create col/linear regression={y=ly}}]
                        {data.dat} node
                        [left,
                         pos=0.875]
                        {\scriptsize$\pgfmathprintnumber{\pgfplotstableregressiona} \cdot x \pgfmathprintnumber[print sign]{\pgfplotstableregressionb}$\quad\ };
            \end{axis}
        \end{tikzpicture}
        }
        \caption{\label{fig:Abb.2}Fehlerbalken}
    \end{figure}
\subsection{GnuPlot}
    \begin{figure}[!ht]
        \centering
        \scalebox{.75}{
        \begin{tikzpicture}
            \begin{axis}[xmin=0, xmax=320,
                         ymin=0, ymax=7.25,
                         axis lines=middle,
                         x post scale = 1.5,
                         y post scale = 1.125,
                         xlabel=$t/s$,
                         ylabel=$\frac{\sigma_a^{NaOH} - \sigma}{\sigma - \sigma_a^{NaAc}}$,
                         x label style={at={(axis description cs:.5,-.1)}, anchor=north},
                         y label style={at={(axis description cs:-.075,.5)}, rotate=90, anchor=south},
                         xmajorgrids=false,
                         ymajorgrids=false,
                         grid style=dashed,
                         legend pos=north west]
                %Messwerte bei 298.15K
                \addplot[only marks,
                         mark = x] table
                        {298.txt};
                %Eichgerade
                \addplot[raw gnuplot,
                         thick,
                         red,
                         dotted,
                         mark=none] gnuplot
                        {f(x)=a*x;
                         a=1;
                         set fit errorvariables;
                         fit f(x) '298.txt' using 1:2 via a;
                         plot [x=0:300] f(x);
                         set print "parameters.dat";
                         print a, a_err;} node
                        [left,
                         pos=0.625]
                        {\pgfplotstableread{parameters.dat}\parameters\pgfplotstablegetelem{0}{0}\of\parameters \pgfmathsetmacro\paramA{\pgfplotsretval}$\frac{\sigma_a^{NaOH} - \sigma}{\sigma - \sigma_a^{NaAc}} = \pgfmathprintnumber{\paramA}\,\frac{1}{s} \cdot t\quad$};
            \end{axis}
        \end{tikzpicture}
        }
        \caption{Ursprungsgerade}
    \end{figure}
    
    \begin{filecontents*}{data.csv}
      X Y
      1.000 0.00
      -2.000 2.00
      3.000 14
    \end{filecontents*}
    \begin{filecontents*}{data2.csv}
      X Y
      -1.000 0.00
      -2.000 2.00
      3.000 -3
    \end{filecontents*}
    
    \begin{figure}[H]
        \centering
        \begin{tikzpicture}
            \begin{axis}[xmin=-3, xmax=4,
                         ymin=0, ymax=15,
                         axis lines=middle,
                         x post scale = 1.5,
                         y post scale = 1.125,
                         xlabel=$t/s$,
                         ylabel=$\frac{\sigma_a^{NaOH} - \sigma}{\sigma - \sigma_a^{NaAc}}$,
                         x label style={at={(axis description cs:.5,-.1)}, anchor=north},
                         y label style={at={(axis description cs:-.075,.5)}, rotate=90, anchor=south},
                         xmajorgrids=false,
                         ymajorgrids=false,
                         grid style=dashed,
                         legend pos=north west]
                \addplot[only marks] table {data.csv};
                \addlegendentry{Data}
            
                \linregAddplot[blue]{data.csv}
                \addlegendentry{$y=\linregSlope \cdot x + \linregIntercept$, ${R}^{2}=\linregRsq$}
            
            \end{axis}
        \end{tikzpicture}
        \caption{Steigung durch vereinfachenden Code in Präambel (s. PGfplotset)}
    \end{figure}
    \begin{figure}[H]
            \centering
            \begin{tikzpicture}
                \begin{axis}[xmin=0, xmax=250,
                             ymin=0, ymax=1.7,
                             axis lines=middle,
                             x post scale = 1.5,
                             y post scale = 1.125,
                             xlabel=$t/s$,
                             ylabel=$\Delta E$,
                             x label style={at={(axis description cs:.5,-.1)}, anchor=north},
                             y label style={at={(axis description cs:-.075,.5)}, rotate=90, anchor=south},
                             xmajorgrids=false,
                             ymajorgrids=false,
                             grid style=dashed,
                             legend pos=north west,
                             legend style={font=\tiny,
                                           cells={anchor=west},
                                           legend columns=2,
                                           %at={(.75,-.225)},
                                           draw=black}]
                    % pH = 8.0 %%%%%%%%%%%%%%%%%%%%
                    \addplot[only marks,
                             MidnightBlue] table
                            {pH80.txt};
                    \addlegendentry{$pH = 8.0$}
                    % pH = 8.0 - Eichgerade
                    \addplot[raw gnuplot,
                             no markers,
                             MidnightBlue] gnuplot
                            {% Calculate a bunch of statistical values
                             stats 'pH80.txt';
                             % plot the linear fit
                             plot [x=0:240] STATS_slope*x + STATS_intercept;
                             % Open a file to save the parameters into
                             set print "values.dat";
                             % Write the parameters to file
                             print STATS_slope;
                             print STATS_intercept;
                             print STATS_correlation;
                             x};
                    % Open the file again
                    \newread\myread
                    \openin\myread=values.dat
                    % Read the values from the file and pprint them to macros
                    \read\myread to \mytmp \pgfmathprintnumberto{\mytmp}{\slope}
                    \read\myread to \mytmp \pgfmathprintnumberto{\mytmp}{\intercept}
                    % Print Legendentry
                    \addlegendentry{$y=\slope \cdot x + \intercept$}
                    
                    % pH = 10.0 %%%%%%%%%%%%%%%%%%%%
                    \addplot[only marks,
                             DarkOrchid] table {pH100.txt};
                    \addlegendentry{$pH = 10.0$}
                    % pH = 10.0 - Eichgerade
                    \addplot[raw gnuplot,
                             no markers,
                             DarkOrchid] gnuplot
                            {stats 'pH100.txt';
                             plot [x=0:240] STATS_slope*x + STATS_intercept;
                             set print "values.dat";
                             print STATS_slope;
                             print STATS_intercept;
                             print STATS_correlation;
                             x};
                    \newread\myread 
                    \openin\myread=values.dat
                    \read\myread to \mytmp \pgfmathprintnumberto{\mytmp}{\slopee} 
                    \read\myread to \mytmp \pgfmathprintnumberto{\mytmp}{\interceptt}
                    \addlegendentry{$y=\slopee \cdot x + \interceptt$}
                \end{axis}
            \end{tikzpicture}
            \caption{Multiple slopes}
        \end{figure}
    
\section{Others}

    \begin{align*}
        3\cdot2
        &= 6\,\si{mm}
        \\[1em]
        &\left. 2x + \frac{2}{3}x^{\frac{3}{2}} + \frac{1}{x}\right|_{1}^{9}
        \\
        &y \cdot x(y)\Bigr|_{y_{1}}^{y_{2}}
        \\
        &2x + \frac{2}{3}x^{\frac{3}{2}} + \frac{1}{x}\Biggr|_{1}^{9}
        \\[1em]
        &\Delta_{\underset{\mathclap{\substack{\uparrow \\ \text{Cambustion}}}}{C}} H^{\circ\kern-0,525em-}_{298K}\qquad(10\,pt)
        \\
        &\Delta_{\underset{\mathclap{\substack{\uparrow \\ \text{Cambustion}}}}{C}} H^{\circ\kern-0,47em-}_{298K}\qquad(12\,pt)
        \\
        &\Delta_{\underset{\mathclap{\substack{\uparrow \\ \text{Cambustion}}}}{C}} H^{\circ\kern-0,65em-}_{298K}
        \\[1em]
        \AA
        &\text{ oder }\si{\angstrom}
        \\[1cm]
        \si{\milli\litre}
        \\
        2\si{ml}
        \\
        \SI{2}{ml}
        \\
        2\,\text{ml}
        \\
        \SI{2}{\milli\l}
    \end{align*}

\end{document}
